\chapter{刑罚的裁量}

% Source PDF physical page 190; see audit/source_coverage.json.
\begin{EditionOpeningPanel}
1.体系内容：量刑情节、累犯、自首、立功、数罪并罚、缓刑。

2.重要考点：一般累犯的成立条件、自首与立功的情形、数罪并罚的运用、缓刑的撤销。
\end{EditionOpeningPanel}

\section{量刑情节}

法定量刑情节有：从重、从轻、减轻、免除处罚等情节。

1.从重处罚和从轻处罚

这是指在同一个刑格幅度内，从重处罚或从轻处罚。

2.减轻处罚

第63条第1款犯罪分子具有本法规定的减轻处罚情节的，应当在法定刑以下判处刑罚；本法规定有数个量刑幅度的，应当在法定量刑幅度的下一个量刑幅度内判处刑罚。

（1）减轻处罚，是指在刑格以下判处刑罚。

[注意]这里的“以下”不包括本数。例如，对盗窃罪的“三年以上十年以下有期徒刑”，如果判3年，属于从轻处罚，而非减轻处罚。从轻处罚和减轻处罚不存在竞合。

（2）减轻处罚有刑格限制，只能在下一个刑格内处罚，而不能跨越下一个刑格，在下下一个刑格内处罚。

（3）特别减轻处罚制度。第63条第2款：“犯罪分子虽然不具有本法规定的减轻处罚情节，但是根据案件的特殊情况，经\CJKunderwave{最高人民法院核准}，也可以在法定刑以下判处刑罚。”（2024年试题）

（4）减轻处罚包括对主刑减轻处罚，也包括对附加刑减轻处罚。\footnote{参见最高人民法院第二巡回法庭2019年第15次法官会议纪要《减轻处罚应同时适用于附加刑》。}（2024年试题）

3.免除处罚

（1）这是指对犯罪分子作出有罪宣告，但是免除其刑罚处罚。例如，第24条第2款规定：对于中止犯，没有造成损害的，应当免除处罚…·

（2）免除处罚在性质上不同于无罪判决或不起诉，免除处罚的前提是有罪判决。

% Source PDF physical page 191; see audit/source_coverage.json.
（3）免除处罚是指免除刑罚处罚，不意味着免除非刑罚处罚，例如，还可判处训诫、责令赔偿损失等非刑罚处罚。如果既免除刑罚处罚，也免除非刑罚处罚，就是单纯宣告有罪。因此，免除处罚不等于单纯宣告有罪。

例如，狗蛋构成交通肇事罪，并且“肇事后逃逸”，应适用的刑格是“3年以上7年以下有期徒刑”。

\InsertDiagram{figures/manual_p191_1.tex}

\section{累犯}

[问题］狗蛋犯故意伤害罪被判3年有期徒刑，刑罚执行完毕后，第二年又犯危险驾驶罪。狗蛋是否构成累犯？（危险驾驶案）

累犯，是指刑罚执行完毕或赦免以后，在一定期间内又犯罪的人。

\subsection*{一、一般累犯}

第65条被判处\CJKunderwave{有期徒刑以上刑罚}的犯罪分子，\CJKunderwave{刑罚执行完毕或者救免以后}，在\CJKunderwave{五}\CJKunderwave{年以内}再犯\CJKunderwave{应当判处有期徒刑以上刑罚}之罪的，是累犯，应当从重处罚，但是\CJKunderwave{过失犯罪}和\CJKunderwave{不满十八周岁的人犯罪}的除外。

前款规定的期限，对于被假释的犯罪分子，从\CJKunderwave{假释期满之日}起计算。

1.主观条件：前后罪都必须是故意犯罪。

2.年龄条件：前后罪都必须是已满18周岁的人犯罪。第一次犯罪时行为人未满18周岁，第二次犯罪时行为人已满18周岁，不是累犯。（2015年第59题，2017年第56题）

3.刑度条件：前后罪都必须是被判或应判有期徒刑以上刑罚的犯罪。例如，本节开头“危险驾驶案”中，由于危险驾驶罪的最高刑是拘役，故不构成累犯。

4.时间条件：后罪发生的时间，必须是在前罪\CJKunderwave{刑罚执行完毕}或者\CJKunderwave{赦免}以后5年内。

（1）前罪的“刑罚”范围。前罪刑罚执行完毕，是指主刑（有期徒刑）执行完毕。

第一，前罪刑罚不包括拘役。例如，甲犯故意伤害罪被判6个月拘役，犯过失致人死亡罪被判2年有期徒刑，数罪并罚，只执行2年有期徒刑，刑罚执行完毕后第3年又犯故意杀人罪，应判无期徒刑。由于前罪中故意犯罪只是判拘役，没有得到执行，甲不能构成累犯。

第二，前罪刑罚不包括管制。管制是否执行完毕，在所不问。例如，甲犯A罪被判3年有期徒刑，犯B罪被判1年管制。2010年6月底有期徒刑执行完毕，2011年2月又犯盗窃罪，应判1年有期徒刑。管制的执行期间是2010年7月至2011年7月。虽然甲在此期间又犯盗窃罪，仍构成累犯。

第三，前罪刑罚不包括附加刑。附加刑是否执行完毕，在所不问。例如，甲犯放火罪

% Source PDF physical page 192; see audit/source_coverage.json.
被判3年有期徒刑，剥夺政治权利2年。2010年6月底有期徒刑执行完毕，2011年6月又犯盗窃罪，应判1年有期徒刑。剥夺政治权利的期间是2010年7月至2012年7月。虽然甲在此期间又犯盗窃罪，仍构成累犯。

（2）前罪刑罚要求“执行完毕”，而不是“不再执行”。

第一，假释的问题。首先，在假释考验期内犯新罪，不成立累犯，因为刑罚没有执行完毕，此时要撤销假释，数罪并罚；其次，在假释期满后犯罪，可以成立累犯，因为成功的假释视为原判刑罚已经执行完毕。此时累犯的5年起算时间，从假释期满之日起计算，而不是从假释之日起计算。

第二，缓刑的问题。首先，在缓刑考验期内犯新罪，不成立累犯，因为刑罚没有执行完毕，此时要撤销缓刑，数罪并罚；其次，在缓刑考验期满后再犯新罪，也不能成立累犯。因为成功的缓刑，视为原判刑罚不再执行，而非视为原判刑罚已经执行完毕。因为不再执行，就不存在执行完毕问题。因为不属于执行完毕，所以不存在累犯，此时所犯新罪应作为单独犯罪处理。

（3）前罪刑罚执行完毕“5年内”的起算日期：从刑满释放之日起计算。\footnote{2018年12月30日最高人民检察院《关于认定累犯如何确定刑罚执行完毕以后“五年以内”起始日期的批复》。}注意，\CJKunderwave{不}\CJKunderwave{是刑满释放的次日}。例如，甲因盗窃罪被判3年有期徒刑，2019年5月1日上午刑满释放，下午又实施抢劫罪。甲构成累犯。

（4）前罪刑罚的结束方式，不仅包括刑罚执行完毕，还包括赦免。例如（2019年试题），甲因故意伤害罪被判2年有期徒刑，经过特救释放出来，第二年又犯盗窃罪，被判3年有期徒刑。甲构成累犯。

5.累犯（包括特殊累犯）的法律后果：（1）应当从重处罚；（2）不能适用缓刑；

（3）不能适用假释；（4）如果是死缓犯，可以限制减刑。（2015年第10题）

\subsection*{二、特殊累犯}

第66条危害国家安全犯罪、恐怖活动犯罪、黑社会性质的组织犯罪的犯罪分子，在刑罚执行完毕或者赦免以后，在任何时候再犯上述任一类罪的，都以累犯论处。

1.前后罪的种类要求

前后罪是\CJKunderwave{危害国家安全犯罪、恐怖活动犯罪、黑社会性质的组织犯罪}这三类罪中任一类罪。也即前罪与后罪只要是这三类罪即可，不要求保持一致。（2015年第10题）

（1）危害国家安全犯罪包括刑法分则第一章“危害国家安全罪”的所有罪名。

（2）恐怖活动犯罪不仅包括组织、领导、参加恐怖组织罪，帮助恐怖活动罪等；而且包括恐怖组织实施的各种犯罪。

（3）黑社会性质的组织犯罪不仅包括组织、领导、参加黑社会性质组织罪，人境发展黑社会组织罪，包庇、纵容黑社会性质组织罪；而且包括黑社会性质组织实施的各种犯罪。

2.前罪“刑罚执行完毕或者赦免以后”的理解。

（1）前罪判处缓刑，缓刑考验期满，不视为刑罚执行完毕，因此此后不构成累犯。

（2）前罪有数罪，数罪中有拘役的问题。例如，甲因盗窃罪被判3年有期徒刑，因犯参加黑社会性质组织罪被判处6个月拘役，数罪并罚执行3年有期徒刑。该刑罚执行完毕

% Source PDF physical page 193; see audit/source_coverage.json.
之后第6年，甲又犯参加恐怖组织罪，被判1年有期徒刑。首先，由于在第6年犯罪，甲不构成一般累犯。其次，甲要构成特殊累犯，要求前罪中的参加黑社会性质组织罪“刑罚执行完毕”，但是其被判的6个月拘役没有被执行，不属于刑罚执行完毕，故不构成特殊累犯。一般说有期徒刑“吸收”了拘役，其实按照第69条第2款，是指只执行有期徒刑，不执行拘役。

3.成立特殊累犯不要求的事项

（1）前后罪不要求是被判或应判有期徒刑以上刑罚的犯罪。判处拘役、管制、\CJKunderwave{单处附}\CJKunderwave{加刑}都可以。（2017年第56题）

（2）后罪发生的时间，不要求是在前罪刑罚执行完毕或者赦免以后5年内。

\InsertDiagram{figures/manual_table_p193_1.tex}

\section{自首}

[问题］狗蛋杀了人后，不去公安局投案，而是坐等警察上门来抓自已，警察赶到，狗蛋很配合。狗蛋是否构成自首？（上门服务案）

第67条第1款[一般自首］犯罪以后\CJKunderwave{自动投案}，\CJKunderwave{如实供述}自己的罪行的，是自首。对于自首的犯罪分子，可以从轻或者减轻处罚。其中，犯罪较轻的，可以免除处罚。

第2款［准自首］被采取强制措施的犯罪嫌疑人、被告人和正在服刑的罪犯，如实供述司法机关还未掌握的本人其他罪行的，以自首论。

第3款［坦白］犯罪嫌疑人虽不具有前两款规定的自首情节，但是如实供述自己罪行的，可以从轻处罚；因其如实供述自己罪行，避免特别严重后果发生的，可以减轻处罚。

\subsection*{一、一般自首}

关于自首的成立条件，最高人民法院有司法解释，\footnote{2010年12月22日最高人民法院《关于处理自首和立功若干具体问题的意见》。}但缺乏条理，不方便记忆。现通过整理排序，归纳如下。

\subsubsection*{（一）自动投案}

1.\CJKunderwave{投案时间}：犯罪成立后一被动归案前。

（1）“犯罪成立后”，不是指犯罪终局（如既遂）后。例如，甲在杀乙过程中，放弃杀乙，并自动投案，如实供述，构成自首，也成立犯罪中止。

（2）“被动归案前”，是指在被讯问或采取强制措施之前，都可以自动投案。

% Source PDF physical page 194; see audit/source_coverage.json.
（3）在此期间，司法机关是否已经掌握犯罪人的罪行，不影响自动投案的时间窗口。

例如，甲杀人后潜逃，司法机关已经掌握了甲的杀人罪行。此后，甲自动投案，如实供述，成立自首。（2021年试题）

2.\CJKunderwave{投案方式}：主动将自己置于司法机关控制之下，就属于自动投案。

（1）半路上“截胡”。经查实，确已准备去投案，或者正在投案途中，被公安机关捕获的，应当视为自动投案。

（2）坐等警察上门服务。第一，犯罪后主动报案，虽未表明自己是作案人，但没有逃离现场，在司法机关询问时交代自己罪行的，视为自动投案。第二，明知他人报案而\CJKunderwave{在现}\CJKunderwave{场等待}，抓捕时无拒捕行为，供认犯罪事实的，视为自动投案。例如，本节开头“上门服务案”中，狗蛋构成自首。

（3）证据证明程度问题。第一，司法机关仅因形迹可疑，进行盘问、教育后，主动交代自己罪行的，视为自动投案（2017年第9题）。第二，司法机关尚未确定犯罪嫌疑人，进行一般性排查、询问时，主动交代自己罪行的，视为自动投案。第三，人赃俱获时，失去自动投案机会。这是指司法机关在盘问、排查时，在现场发现与犯罪有关的物品，此时交代自己罪行，不能视为自动投案，只能认定为坦白。

（4）在行政拘留、司法拘留、强制隔离戒毒等行政、司法强制措施期间，主动交代司法机关尚未掌握的罪行，视为自动投案。因为如果不主动交代，这些措施届满，行为人就会被释放。行为人主动交代罪行，属于主动将自已置于司法机关刑事措施控制之下。例如，甲因被治安拘留，在拘留期间如实供述自己的强奸事实，成立自首。\footnote{注意：这种自首属于一般自首，而非特别自首。由于行政拘留等措施针对的是违法行为，不是犯罪行为，所以只要行为人交代司法机关尚未掌握的罪行即可，不存在也不要求交代不同种的其他罪行。}

（5）送首。这是指，未成年人或亲属犯罪后，由监护人或亲友送到司法机关。送子归案就属于此种情形。注意：如果犯罪嫌疑人明显反抗，亲友被迫采取捆绑等手段送至司法机关，则不属于自动投案。

3.\CJKunderwave{投案对象}：包括司法机关、非司法机关（例如，犯罪人所在单位、基层组织）、有关个人（例如，单位负责人、被害人）。

[提示】告诉才处理的案件，向被害人主动承认自己的犯罪事实，并自愿归案，视为自首。例如，甲暗中诽谤乙，然后向乙主动承认是自己所为，并自愿归案，构成自首。

4.\CJKunderwave{投案彻底性}：投案后自愿接受控制，直到最终审判。例如，甲犯罪后，先向警方打电话如实供述，但随后不归案，不成立自首。注意：逃跑又回来的两种情形。（2022年试题）

情形一：甲犯A罪后，自动投案，如实供述，成立自首。此后甲从看守所逃跑。对甲撤销自首。并且逃跑行为构成脱逃罪。此后甲又自动投案。一方面恢复A罪的自首，另一方面构成脱逃罪的自首，想象竞合。（2010年主观题）

情形二：甲犯A罪后，被动归案，不成立自首。此后甲从看守所逃跑，构成脱逃罪。此后甲又自动投案。甲的这个“自动投案”不构成A罪的自首，因为此前因A罪被动归案后，就失去针对A罪的自首机会：就A罪而言，甲的这个“自动投案”没有节约司法资源，只是恢复到原来被动归案的状态。不过，甲的这个“自动投案”构成脱逃罪的自首。（2009年主观题，2022年试题）

% Source PDF physical page 195; see audit/source_coverage.json.
\subsubsection*{（二）如实供述}

1.如实

（1）如实的标准。主观标准认为（主流观点），如实供述是指诚实供述，也即犯罪人诚实地将自己的记忆内容供述出来，即使与事后查明的客观犯罪事实不一致，也算如实供述。客观标准认为，成立如实供述，要求犯罪人供述内容必须与事后查明的客观事实相符。

（2）关于法律适用的辩解不影响如实供述。如实供述了案件事实，但对案件事实的定性存在不同理解，有不同看法，进行辩解，仍属于如实供述。也即，只要求就“事实判断”如实供述，至于“价值评价（法律适用）”上不作要求。

例1（2009年第53题），甲自动投案，如实供述了伤害行为，但一直主张自已是正当防卫，不影响如实供述。

例2（2017年第9题），甲自动投案，如实供述了杀人事实，司法机关认定其主观罪过是故意，甲辩称是过失，不影响如实供述。

例3，认为自已不是贪污，只是经济问题，不影响如实供述的成立。

2.供述

（1）供述内容：第一，只要求供述影响定罪量刑的犯罪事实，不要求供述所有细节。

例如（2009年第53题），甲自动投案后，如实供述了杀人行为，但拒绝说明凶器藏匿地点，视为如实供述，因为“藏匿凶器”不影响定罪量刑。

第二，供述内容还包括姓名、年龄、职业、住址、前科等基本情况。这些信息若不影响定罪量刑，则不影响如实供述。例如，隐瞒自已国家工作人员身份，逃避贪污罪的追究，不属于如实供述。（2017年第9题）

第三，对如实供述的内容可以分别认定自首。例如，甲如实供述了自己抢劫罪的事实，但隐瞒了自己持枪抢劫的加重处罚情节。就抢劫罪而言，可以成立自首。

（2）截止时间。自动投案，如实供述，成立自首；后又翻供的，撤销自首；但在一审判决前又如实供述，恢复自首。（2009年第53题）

\subsection*{二、特别自首（准自首、余罪自首）}

这是指，甲因A罪行被采取强制措施后，甲如实供述司法机关尚未掌握的B罪行。

1.主体：必须是依法被采取强制措施的犯罪嫌疑人、被告人和正在服刑的罪犯。

[提示］这里的强制措施包括监察机关为调查犯罪而采取的留置措施。

2.“B罪行”

（1）“B罪行”必须是司法机关尚未掌握的罪行。

（2）“B罪行”与“A罪行”必须是不同种罪行。根据司法解释，判断标准是：一般而言，罪名不同，便是不同种罪行。例外：虽然罪名不同，但二者具有密切联系，仍属于同种罪行。密切联系的表现有：二者属于选择性罪名，二者具有吸收关系、牵连关系。

例1，甲因受贿被捕，又交代因受贿为他人谋取利益的行为，构成滥用职权罪、徇私枉法罪的，应认定为同种罪行。

例2，甲因出售假币被捕，又交代自己购买假币的罪行，应认定为同种罪行。

例3，甲因组织卖淫被捕，又交代自已强迫卖淫的罪行，应认定为同种罪行。

% Source PDF physical page 196; see audit/source_coverage.json.
[注意］如果经查证A罪行不成立，则如实供述司法机关尚未掌握的罪行，无论与A罪行是不是同种罪行，都成立准自首。

例1（2009年第53题），甲涉嫌贪污罪被采取强制措施后，司法机关发现贪污罪不成立。此时，甲又主动供述司法机关尚未掌握的贪污罪事实。甲构成准自首。

例2，甲涉嫌受贿罪被采取强制措施后，司法机关发现受贿罪不成立。此时，甲又主动供述司法机关尚未掌握的滥用职权罪的事实。甲构成准自首。

\subsection*{三、坦白}

这是指，甲因A罪行被采取强制措施后，甲如实供述司法机关已经掌握的A罪行。坦白是法定从宽处罚情节，但是从宽力度比自首小。

1.坦白和一般自首的区别在于：坦白是被动归案，一般自首是自动投案。

2.坦白和特别自首的相同点在于：二者都是被动归案。区别在于：坦白是如实供述司法机关已经掌握的罪行，特别自首是如实供述司法机关尚未掌握的罪行。

\subsection*{四、自首的特别问题}

1.数罪自首：犯有数罪的犯罪嫌疑人仅如实供述所犯数罪中部分犯罪的，对如实供述的部分犯罪，认定为自首。也即，各算各的账。

2.交通肇事罪的自首：逃逸后又归案，属于自首；没有逃逸，直接主动报案、归案，如实供述，也是自首。（2017年第9题）

3.单位犯罪的自首。总结司法解释的要点：

第一，单位可以成为自首的主体。例如，单位集体决定自首，然后派代表去自首。

第二，单位没有自首，单位的主管人员（领导）自首，视为单位自首。

第三，单位没有自首，单位的直接责任人员（直接办事人）自首，只能认定为其个人自首，不能视为单位自首。

第四，单位自首了，单位的主管人员和直接责任人员的坦白，视为自首。

[引申总结]单位犯罪的犯罪中止。

第一，单位没有犯罪中止，单位的主管人员（领导）犯罪中止，视为单位犯罪中止。

第二，单位没有犯罪中止，单位的直接责任人员（直接办事人）犯罪中止的，只能认定为其个人犯罪中止，不能视为单位犯罪中止。并且，由此导致单位犯罪未能得逞，单位犯罪构成犯罪未遂。（2012年第54题）

4.一般自首与犯罪中止的对比。

（1）二者相同点：都具有自动性，都是法定从宽处罚事由。

（2）二者不同点：第一，自动性的内容不同。自首是自动投案。犯罪中止是自动放弃犯罪。例如，甲杀害乙，捅了两刀，然后打电话报案，坐等警察来抓，等待期间，发现乙没死，又杀死乙，然后继续等待，警察到后，配合抓捕，并如实供述罪行。甲杀死乙后，继续等待抓捕，并如实供述，成立自首。有法官认为甲不成立自首，理由是甲报案后还继续杀人，说明没有自动放弃犯罪。然而，这种看法将自首的自动性的内容与犯罪中止的自动性的内容混淆了，自首只要求自动投案。

% Source PDF physical page 197; see audit/source_coverage.json.
第二，自动性的时间不同。自首是犯罪实施后，被动归案前。犯罪中止是犯罪实施后，终局形态形成前。二者的起始时间相同，截止时间不同。注意，自首的起始时间，是犯罪实施以后，而不是犯罪终局以后，在犯罪过程中可以自首。例如，甲在杀乙过程中，自动放弃杀乙，并自动投案，如实供述，构成自首。（2017年第9题）

（3）二者关系。虽然二者有不同点，但二者不是“A与非A”的对立排斥关系，而是“A与B”的两码事关系（中立关系），可以并存。例如，甲在杀乙过程中，自动放弃杀乙，并将乙抢救成功，同时自动投案，如实供述，构成自首，也构成犯罪中止。

\section{立功}

[问题]狗蛋因盗窃罪被抓后，主动向警方提供同案犯狗剩的姓名、体貌特征。狗剩因此被抓。狗蛋是否构成立功？（投名状案）

\subsection*{一、立功的方式、根据、时间}

\subsubsection*{（一）立功的方式}

立功是法定从宽处罚情节。根据刑法第68、78条规定，立功的方式有五种：（1）揭发他人犯罪；（2）协助抓捕犯罪人；（3）阻止他人犯罪；（4）提供犯罪线索；（5）对国家和社会有重大贡献。

\subsubsection*{（二）立功的根据}

第1-4项立功，从宽处罚的根据是，节约了侦查机关的司法资源，也即帮侦查机关省事了，予以奖励。

第5项立功，从宽处罚的根据是，有利于国家和社会。对国家和社会有重大贡献，主要包括：（1）有发明专利；（2）舍已救人；（3）在抢险抗灾中有突出表现。注意：发明专利问题。第一，不包括实用新型专利和外观设计专利。第二，必须是犯罪人独立完成或者为主完成。

\subsubsection*{（三）立功的时间}

1.第1-4项立功的时间：到案后。\footnote{少数说认为是，犯罪后。例如，甲乙丙共同杀人后，各自逃匿。过了一周，甲向公安机关打电话：“我和乙丙杀了人，如果我先投案，乙丙会杀我家人。我告诉你们，乙丙犯罪后的藏匿地点，你们抓到乙丙，我再投案。”公安机关依照甲提供的地址，抓到乙丙。然后甲投案自首。第一，甲构成自首。第二，甲提供地址，协助抓捕乙丙，是否构成立功？多数说认为不构成，因为甲没有到案。少数说认为构成，因为甲是在犯罪后协助抓捕乙丙。}

提示：犯罪分子到案后，又逃跑，在逃跑期间，可构成立功。也即，只要第一次到案后，此后就可以立功。例如，甲犯罪后被动归案，又逃跑，在逃跑期间，甲给警察打电话，揭发他人犯罪。甲构成立功。（2022年试题）

2.第5项立功的时间：到案后，包括服刑期间。

% Source PDF physical page 198; see audit/source_coverage.json.
\subsection*{二、揭发他人犯罪}

\subsubsection*{（一）揭发他人“\CJKunderwave{犯罪}"\protect\footnotemark}\footnotetext{参见张明楷：《论揭发他人“犯罪行为”》，载《人民检察》2006年第15期。}

1.这里的他人“犯罪”，要求是“制造违法事实”意义上的犯罪。

例1，犯罪人甲揭发乙犯故意伤害罪。经查，乙构成正当防卫或紧急避险。正当防卫和紧急避险是排除违法事由。这表明，乙的行为不是“制造违法事实”意义上的“犯罪”行为。甲不构成立功。（2023年试题）

例2，犯罪人甲揭发乙实施了强奸罪。经查，乙有强奸事实，但当时只有13岁。乙的行为是“制造违法事实”意义上的“犯罪”行为。甲属于立功。

例3，犯罪人甲揭发乙将丙打成轻微伤。由于成立故意伤害罪要求达到轻伤，所以乙的行为不构成犯罪。乙的行为不是“制造违法事实”意义上的“犯罪”行为。甲不是在揭发他人“犯罪”，不构成立功。（2023年试题）

2.揭发行为要节约了侦查机关的司法资源，才成立立功。

例1，犯罪人甲揭发乙犯告诉才处理的犯罪。告诉才处理的案件，是原告到法院去起诉，不需要公安机关立案侦查，因此，不存在节约公安机关司法资源的问题。所以，甲不构成立功。

例2，甲揭发乙犯抢劫罪。对此，公安机关须立案侦查。经查，乙有抢劫事实，但侦破时乙已经死亡或已过追诉时效。由于甲提供犯罪线索，为公安机关省事了，所以，甲构成立功。

\subsubsection*{（二）“\CJKunderwave{揭发}”他人犯罪}

这里的揭发包括犯罪人告发他人对自已犯罪。例如，甲女因诈骗罪被捕，揭发乙男强奸自己，应认定为立功。换言之，这里的揭发包含被害人告发，否则对被害人不公平，因为被害人也节约了司法资源。

\subsubsection*{（三）揭发“\CJKunderwave{他人}”犯罪}

1.这里的“他人”不包括共犯人。\CJKunderwave{自首是供述自己的罪行，立功是揭发他人的罪行}。揭发共犯人的共同犯罪，属于“如实供述自己的罪行”，是成立自首、坦白的条件，不认定为立功（因为没有额外地节约司法资源）。例如，甲乙共同盗窃。甲自动投案，如实供述自己和乙共同盗窃。甲构成自首，不构成立功。如果甲隐瞒了同案犯乙，谎称自己一个人盗窃，则不构成自首。

2.揭发共犯人共同犯罪以外的其他犯罪，构成立功。例如，甲乙共同盗窃。甲自动投案，如实供述了自已和乙共同盗窃。甲构成自首。并且，甲揭发乙还单独实施过绑架罪。甲构成立功。

[总结］通过“出卖兄弟”（同案犯），能够获得三项好处：（1）成立自首；（2）揭发犯罪型立功；（3）协助抓捕型立功，也即协助抓捕同案犯。（2024年主观题）

3.对向犯的问题

（1）处罚双方的对向犯中，揭发相对方的犯罪行为，只构成自首（或坦白），不构成

% Source PDF physical page 199; see audit/source_coverage.json.
立功。这是因为，双方在事实层面是共同犯罪，\footnote{最终处理方案，不一定按照共同犯罪处理。}揭发相对方的共同犯罪行为，是如实供述自己罪行的必要组成部分，没有额外地节约司法资源。（新变化的观点）

例1，甲向乙购买了枪支。后来甲自动投案，如实供述自已购买枪支的事实，并交代了是从乙处购买的。甲只有将卖家乙供出来，才属于“如实供述”了自已的非法买卖枪支罪，才成立自首。因此，甲不构成立功。

例2，甲向官员乙行贿，乙接受贿赂。后来乙自动投案，如实供述自己受贿，并交代了甲行贿。乙交代甲行贿，是如实供述自已受贿罪的必要组成部分，构成自首，不构成立功。\footnote{参见2013年1月1日最高人民法院、最高人民检察院《关于办理行贿刑事案件具体应用法律若干问题的解释》第7条、第9条规定。}

[注意］揭发相对方的共同犯罪行为，如果不认定为自首，则可以认定为立功。在此，立功与自首是补充关系，也即优先认定为自首，若不能认定为自首，则可以考虑认定为立功，因为节约了司法资源，要奖励，不奖励为自首，就要奖励为立功。

例如，甲的亲属犯罪，甲送给公安局长乙10万元，乙收钱，利用职权使甲的亲属无罪。乙的该行为构成徇私枉法罪。乙因徇私枉法罪被捕，乙如实供述了甲行贿、自已受贿的事实。根据司法解释，由于乙的受贿罪与徇私枉法罪具有密切联系，属于同种罪行，因此，乙主动交代自己的受贿罪，不构成特别自首。但是，乙也揭发了甲的行贿罪，就此而言，构成立功。（2023年主观题）

（2）处罚一方的对向犯中，揭发相对方的犯罪行为，构成立功，因为节约了司法资源。

例1，甲因诈骗罪被抓，主动交代自己从乙处购买淫秽物品。甲属于揭发乙犯贩卖淫秽物品牟利罪，构成立功。

例2，甲因诈骗罪被抓，主动交代，自己为了吸食，从乙处购买毒品。甲属于揭发乙犯贩卖毒品罪，构成立功。

（3）在如实供述自己罪行之外，揭发其他犯罪的，可构成立功，因为节约了司法资源。

例1，甲在家里放了许多毒品，用于自已吸食。后甲自动投案，如实供述非法持有毒品的事实，并且交代了是从乙处购买的。甲只要如实供述了非法持有毒品的事实，不用交代卖家，就构成非法持有毒品罪的自首。甲还交代了卖家乙，揭发乙的贩卖毒品罪。甲揭发乙，属于在自首之外额外多干了一件事，构成立功。甲既享受自首待遇，还享受立功待遇。

例2，甲贩卖毒品。甲被抓后，如实供述贩卖毒品的事实，并且交代了货是从上家乙处购买的。甲要就贩卖毒品罪成立坦白，只需要如实供述贩卖毒品的事实即可，不需要将上家交代出来。因此，将上家交代出来，属于额外多干了一件事，揭发他人犯罪，成立立功。（2019年、2023年、2025年试题）

\subsection*{三、协助抓捕犯罪人}

犯罪分子甲具有下列行为之一，使司法机关抓获犯罪嫌疑人乙的，构成立功：

% Source PDF physical page 200; see audit/source_coverage.json.
1.按照司法机关的安排，以打电话、发信息等方式将犯罪嫌疑人（\CJKunderwave{包括同案犯}）约至指定地点。（2012年第57题）

2.按照司法机关的安排，当场指认、辨认犯罪嫌疑人（\CJKunderwave{包括同案犯}）。

3.带领侦查人员抓获犯罪嫌疑人（\CJKunderwave{包括同案犯}）。

4.提供司法机关尚未掌握的\CJKunderwave{其他案件的犯罪嫌疑人}的联络方式、藏匿地址。

5.犯罪分子提供同案犯犯罪后的新信息，例如，新的联络方式、藏匿地址，才可能构成立功。\CJKunderwave{注意}：提供同案犯的姓名、住址、体貌特征等基本情况，或者提供犯罪前、犯罪中掌握的同案犯联络方式、藏匿地址，这些信息都不属于犯罪后的新信息，不构成立功。例如，本节开头“投名状案”中，狗蛋不构成立功。（2012年第57题）

\subsection*{四、提供犯罪线索}

“提供犯罪线索”型立功，与“揭发犯罪”型立功、“协助抓捕”型立功、“阻止犯罪”型立功，不是对立排斥关系，是后三者的方式之一。例如，通过提供他人犯罪的线索，可以揭发他人犯罪、协助抓捕他人、阻止他人犯罪等。

不过，提供犯罪线索，还有其他一些效果。例如，提供重要证据、证人，使案件得以侦破。

该犯罪线索应是司法机关尚未掌握的犯罪线索，并且查证属实。司法解释要点：\footnote{2010年12月22日最高人民法院《关于处理自首和立功若干具体问题的意见》。}

1.犯罪分子通过暴力、胁迫、贿买等非法手段，获取他人犯罪线索，“检举揭发”的，不构成立功。

（1）这里的“贿买”，是指向国家工作人员行贿而购买。向非国家工作人员购买或免费获取（无欺诈、胁迫等非法手段）的犯罪线索，检举揭发的，构成立功。这里的非国家工作人员包括同监犯。也即，同监犯将自己或他人尚未被司法机关掌握的犯罪线索，卖给或无偿提供给犯罪分子，犯罪分子检举揭发的，构成立功。

（2）犯罪分子从负有查办犯罪、监管职责的国家工作人员处获取他人犯罪线索，无论是通过贿买而获取，还是免费获取，检举揭发的，都不构成立功。这是因为，提供这种线索，没有节约司法资源。

（3）通过非法手段（如引诱、欺骗、胁迫）制造他人犯罪，获取他人犯罪线索，检举揭发的，不构成立功。（2023年试题）

2.总结亲友问题

（1）犯罪分子的亲友亲自向司法机关提供他人犯罪线索，不能认定为犯罪分子立功。（2012年第57题）

（2）犯罪分子被羁押后，律师、亲友在会见过程中，违反监管规定，将他人犯罪线索告知犯罪分子，犯罪分子“检举揭发”的，不构成立功。

（3）犯罪分子被羁押前，亲友将他人犯罪线索告知犯罪分子，犯罪分子检举揭发的，构成立功。亲友获得犯罪线索的方式可以是购买，但不能是非法手段。（2022年试题）

3.犯罪分子将本人以往查办犯罪职务活动中掌握的他人犯罪线索予以检举揭发的，不构成立功。

% Source PDF physical page 201; see audit/source_coverage.json.
[总结］一般自首、特别自首、坦白、立功。

\InsertDiagram{figures/manual_table_p201_1.tex}

[注意］自首、坦白、立功的法律效果

例如，甲犯了A罪、B罪。甲自动投案，如实供述了A罪罪行，并且揭发乙犯C罪。甲构成自首，从宽处罚的效果只能惠及于A罪（若甲被动归案，如实供述A罪罪行，构成坦白，则坦白的从宽处罚效果也只能惠及于A罪）。甲构成立功，从宽处罚的效果能够惠及于A罪和B罪。（2025年试题）

\section{数罪并罚}

［问题］狗蛋犯甲罪，被判有期徒刑10年；犯乙罪，被判有期徒刑13年；犯丙罪，被判有期徒刑15年；犯丁罪，被判管制1年。对狗蛋如何数罪并罚？（数罪并罚案）

数罪并罚，是指法院对犯罪人所犯数罪分别定罪量刑后，按照法定数罪并罚原则决定应执行的刑罚。

\subsection*{一、数罪并罚的原则}

第69条判决宣告以前一人犯数罪的，除判处死刑和无期徒刑的以外，应当\CJKunderwave{在总和}\CJKunderwave{刑期以下、数刑中最高刑期以上}，酌情决定执行的刑期，但是管制最高不能超过三年，拘役最高不能超过一年，有期徒刑总和刑期不满三十五年的，最高不能超过二十年，总和刑期在三十五年以上的，最高不能超过二十五年。

\CJKunderwave{数罪中有判处有期徒刑和拘役的，执行有期徒刑。数罪中有判处有期徒刑和管制，或}\CJKunderwave{者拘役和管制的，有期徒刑、拘役执行完毕后，管制仍须执行}。[本款由《刑法修正案

\subsubsection*{（九）》所增设]}

数罪中有判处附加刑的，附加刑仍须执行，其中附加刑种类相同的，合并执行，种类不同的，分别执行。

1．吸收原则

吸收了某个刑的意思是\CJKunderwave{不执行这个刑}，而不是将其吸收进来一并执行。三种吸收情形：

（1）死刑（作为最高刑）+其他主刑=死刑（不执行其他主刑）。

（2）无期（作为最高刑）+其他主刑=无期（不执行其他主刑）。

（3）有期（作为最高刑）+拘役=有期（不执行拘役）。（2016年第55题，2017年第55题）

% Source PDF physical page 202; see audit/source_coverage.json.
2.限制加重原则（不是无限累加，而是有个上限）

（1）死刑+死刑=死刑（只执行一个死刑）。

（2）无期+无期=无期（只执行一个无期）。

（3）有期徒刑+有期徒刑。单罪有期徒刑是6个月至15年。数罪并罚分三种情形。

例如，狗蛋犯三个罪。情形一：分别判处2年，5年，10年。情形二：分别判处10年，12年，11年。情形三：分别判处10年，12年，14年。（2012年第12题）

\InsertDiagram{figures/manual_p202_1.tex}

情形一：下限是10年，上限是17年。

情形二：下限是12年，上限是20年。

情形三：下限是14年，上限是25年。

（4）拘役+拘役。单罪拘役是1个月至6个月。数个拘役，数罪并罚，可以超过6个月达到1年（1-6-1）。

（5）管制+管制。单罪管制是3个月至2年。数个管制，数罪并罚，可以超过2年达到3年（3-2-3)。

3.并科原则（分别执行）

（1）有期（作为最高刑）+管制=有期+管制（分别执行，管制在有期徒刑执行完毕后执行）。例如，本节开头“数罪并罚案”中，三个有期徒刑，下限是15年，上限是25年，比如选择25年，执行完毕后再执行管制1年。

（2）拘役（作为最高刑）+管制=拘役+管制（分别执行，管制在拘役执行完毕后执行）。（2017年第55题）

4.附加刑的并罚问题

（1）数罪中有判处附加刑的，附加刑仍须执行。也即，附加刑不能被主刑吸收。例如，甲犯A罪被判处拘役3个月，附加剥夺政治权利1年，犯B罪被判处有期徒刑2年，附加剥夺政治权利3年。数罪并罚时，就主刑而言，有期徒刑吸收了拘役，只执行有期徒刑2年。就附加刑而言，剥夺政治权利合并执行4年。这个剥夺政治权利不能被有期徒刑吸收，必须执行。

（2）附加刑种类相同的，合并执行。合并执行不是吸收执行，而是累加执行。例如，判处A罪罚金5万元，判处B罪罚金10万元，合并执行就是15万元。

（3）附加刑种类不同的，分别执行。\CJKunderwave{注意}：罚金与没收财产种类不同，应分别执行。对此详见上一讲“没收财产”的总结。

\subsection*{二、发现漏罪，先并后减}

第70条判决宣告以后，刑罚执行完毕以前，发现被判刑的犯罪分子在判决宣告以前还有其他罪没有判决的，应当对新发现的罪作出判决，把前后两个判决所判处的刑罚，依照本法第六十九条的规定，决定执行的刑罚。已经执行的刑期，应当计算在新判决决定的刑期以内。

% Source PDF physical page 203; see audit/source_coverage.json.
例如，甲犯了A罪，判处有期徒刑5年，执行了3年后，发现在判决宣告前，甲还犯了B罪，应判处有期徒刑8年。第一步，先并：将A罪的5年和B罪的8年并，下限是8年，上限是13年。在此期间选择一个刑期比如10年。第二步，后减：10年-已经执行的3年=7年。还需执行7年。总共实际执行10年。

1.当已判决的是两个罪，漏罪跟谁并？结论：跟前两罪决定执行的刑期并。

例如，甲犯有A、B两个罪，A罪判8年，B罪判12年，决定合并执行18年。执行5年后发现判决宣告前，甲还犯有C罪，应判7年。第一步，先并：C罪的7年应与A、B罪的执行刑18年并罚，而不是与A罪的8年、B罪的12年并。下限是18年，上限是20年。比如，选择20年。第二步，后减：20年-已经执行的5年=15年。还需执行15年。总共执行20年。

2.对漏罪判处拘役、管制。

例1（拘役），甲犯某罪被判6年有期徒刑，执行4年后发现漏罪，应被判拘役6个月。先并，6年有期徒刑吸收拘役，只执行6年有期徒刑；后减，6年减去已经执行的4年，还需执行2年。可见，在此对漏罪判处拘役，其实没影响。（2017年第55题）

例2（管制），甲犯某罪被判6年有期徒刑，执行4年后发现漏罪，应被判管制1年。先并，6年有期徒刑与1年管制应分别执行；后减，6年减去4年，还需执行2年有期徒刑，执行完毕以后，还需执行1年管制。

\subsection*{三、又犯新罪，先减后并}

第71条判决宣告以后，刑罚执行完毕以前，被判刑的犯罪分子又犯罪的，应当对新犯的罪作出判决，把前罪没有执行的刑罚和后罪所判处的刑罚，依照本法第六十九条的规定，决定执行的刑罚。

例如，甲犯了A罪，判处有期徒刑5年，执行了3年后，甲又犯B罪，应判处有期徒刑8年。第一步，先减：5年-3年=2年。第二步，后并：2年与8年并。下限是8年，上限是10年。比如，选择10年。还需执行10年。总共实际执行13年。比较发现，“先减后并”比“先并后减”处罚更严厉。

1.对新罪判处拘役、管制。

例1（拘役），甲犯A罪被判5年有期徒刑，执行3年后，又犯故意伤害罪。第一种情形，当时发现了，对故意伤害罪判处6个月拘役。先减，5-3=2；后并，2年有期徒刑吸收拘役，还需执行2年有期徒刑。第二种情形，当时没发现故意伤害罪，A罪的5年执行完毕以后才发现，判6个月拘役。先减，5-5=0；后并，等于再执行6个月拘役。

例2（管制），甲犯某罪判6年有期徒刑，执行4年后又犯新罪，被判管制1年。先减，6-4=2；后并，2年有期徒刑与1年管制应分别执行，也即执行2年有期徒刑完毕以后，还需执行1年管制。

2.倒挂情形。

例如，甲犯A、B、C罪，分别被判13年、10年、15年，合并执行25年；在执行2年后，又犯新罪D罪，应被判8年。对此，先减，25-2=23；后并，23年与8年并罚，此时虽然总和刑期未满35年，但不能因此认为并罚后的最高刑不能超过20年。因为数罪并罚的下限是数刑中最高刑，上限是20年或25年。既然此时下限已是23年，那么决定执行的刑期就应是23年，也即还需执行的刑期。加上已执行的2年，实际执行刑期是25

% Source PDF physical page 204; see audit/source_coverage.json.
年。如果决定执行的刑期是20年，加上已执行的2年，实际执行刑期就只有22年，就会导致又犯新罪，反倒刑期变短。

3.既犯新罪，又发现漏罪：先解决漏罪，再解决新罪，实际操作就是先并后减再并。

例如，甲因盗窃判7年，执行3年后，又犯故意伤害罪应判10年，同时发现漏罪强奸罪应判11年。对此，第一步，先并，7年和11年，假如定15年。第二步，后减，减去3年，剩12年。第三步，再并，再并10年，在12年以上20年以下决定执行刑期。

[总结］总共实际执行刑期的最大值：

\InsertDiagram{figures/manual_table_p204_1.tex}
\SourceFootnote{\textbf{表注①}\quad 例1，甲犯A、B、C罪，分别被判13年、10年、15年，合并执行25年；在执行5年后，发现漏罪D罪，应被判15年。对此，先并，25年并15年，取25年；后减，25年减去已执行的5年，还需执行20年。加上已经执行的5年，实际执行刑期是25年。可见，实际执行刑期不会超过25年。产生这种结果的原因是：对漏罪采取先并后减，“先并”就要受最高刑期的限制，并上去的刑期再多，也不能超过25年。然后再减去已经执行的刑期，即25年-已执行刑期=剩余需执行的刑期。反过来，剩余需执行的刑期+已执行刑期=实际执行刑期，也就是25年。因此，先并后减的实际执行刑期不会超过25年。}
\SourceFootnote{\textbf{表注②}\quad 例2，甲犯A、B、C罪，分别被判13年、10年、15年，合并执行25年；在执行5年后，又犯新罪D罪，应被判15年。对此，先减，25年减去5年是20年；后并，20年并15年，可以决定执行25年，也即还需执行25年。加上已经执行的5年，就是30年。这样实际执行刑期就超过了25年。产生这种结果的原因是：对新罪采取先减后并，“后并”就可能会并到最高刑期25年，而这是剩余仍需执行的刑期。再加上已经执行的刑期，就会超过25年。即：已经执行刑期+剩余仍需执行刑期（25年）=实际执行刑期。}

\section{缓刑}

[问题］狗蛋犯盗窃罪，被判有期徒刑3年，刑满释放后第2年，又犯故意伤害罪，被判有期徒刑1年。对狗蛋能否适用缓刑？（累犯与缓刑案）

\subsection*{一、基本要点}

1.性质：缓刑是暂缓执行刑罚，也即附条件地不再执行刑罚，而非刑罚执行完毕。这与假释的性质不同，假释期满视为刑罚执行完毕。

2.对象条件：被判处3年以下有期徒刑或者拘役的罪犯。

（1）这里的3年以下包括3年。如果是数罪并罚决定执行3年以下有期徒刑，也可以适用缓刑。（2017年第56题，2025年试题）

(2）对被判处管制、单处附加刑的罪犯，不适用缓刑，因为没必要。（2015年第59题)

（3）对于\CJKunderwave{累犯和犯罪集团的首要分子}，不适用缓刑。例如，本节开头“累犯与缓刑案”中，狗蛋又犯故意伤害罪时，构成累犯，因此不能适用缓刑。（2011年第10题，2015年第10题）

3.实质条件：不再有人身危险性。具体要求同时符合下列条件：（1）犯罪情节较轻；

（2）有悔罪表现；（3）没有再犯罪的危险；（4）宣告缓刑对所居住社区没有重大不良影响。

% Source PDF physical page 205; see audit/source_coverage.json.
4.优待对象。一般犯罪人，同时满足上述对象条件和实质条件，只是可以宣告缓刑；而不满18周岁的人、怀孕的妇女、已满75周岁的人，这三类人只要同时满足上述对象条件和实质条件，应当宣告缓刑。

[注意1]这三类人如果不满足上述对象条件和实质条件，则不应宣告缓刑。例如，孕妇和已满75周岁的人若是累犯，就不能宣告缓刑。

[注意2］是否宣告缓刑，与犯罪分子触犯的罪名的性质没有关系。例如，犯罪分子触犯危害国家安全罪，只要被判处3年以下有期徒刑，同时符合实质条件，就可以适用缓刑。又如，犯罪分子触犯故意杀人、爆炸、抢劫、强奸、绑架等暴力性犯罪，只要被判处3年以下有期徒刑，同时符合实质条件，就可以适用缓刑。

5.对宣告缓刑\CJKunderwave{可以}同时适用禁止令，禁止犯罪分子在执行缓刑期间从事特定活动，进入特定区域、场所，接触特定的人。（2011年第10题）

6.在缓刑考验期\CJKunderwave{必须}实行社区矫正。

7.缓刑考验期限：（1）从判决确定之日起计算；（2）判决确定以前先行羁押的日期，不能折抵缓刑考验期限。

8.缓刑只适用于主刑。判处附加刑的，仍需执行附加刑。

\subsection*{二、缓刑的法律后果}

成功的缓刑：原判刑罚不再执行。注意：这不同于原判刑罚已经执行完毕。

失败的缓刑：是指缓刑被撤销。撤销缓刑的三种情形以及处理方式：

1.发现漏罪。这是指发现判决宣告前还有未判决的罪行。

（1）发现的漏罪，是故意犯罪，还是过失犯罪，在所不问；是同种罪，还是异种罪，也在所不问。

（2）撤销缓刑后，对前罪和漏罪进行数罪并罚，然后决定执行的刑期。

[注意】由于已经经过的考验期不算已经执行的刑期，所以这里的数罪并罚不存在“减去已经执行的刑期”的问题，也即不存在先并后“减”。

（3）此后，如果符合缓刑其他条件，可以再次适用缓刑。

（4）缓刑考验期满后才发现漏罪，不能撤销缓刑，只能对漏罪另行起诉审判。

2.又犯新罪。这是指在缓刑考验期内又犯新罪。

（1）又犯的新罪，是故意犯罪，还是过失犯罪，在所不问；是同种罪，还是异种罪，也在所不问。

（2）在缓刑考验期内又犯新罪，无论是在考验期内还是在考验期满后被发现，都要撤销缓刑。\footnote{这是因为，缓刑的实质条件是犯罪分子不再有人身危险性和再犯可能性。如果在考验期内犯新罪，表明有人身危险性和再犯可能性，不符合缓刑条件，应当撤销。因此，即使新罪已经过了追诉时效，也要撤销缓刑，执行原判刑罚，只是不追诉新罪。}

[注意］在数罪并罚后，即使决定执行的刑罚是3年以下有期徒刑，也不能再次适用缓刑。\footnote{这是因为，犯罪人在考验期内又犯新罪，表明其社会危害性和人身危险性仍非常大，就不符合适用缓刑的实质条件。}（2011年第10题）

% Source PDF physical page 206; see audit/source_coverage.json.
（3）在缓刑考验期满后犯新罪，不撤销缓刑。\CJKunderwave{注意}：也不可能构成累犯，因为缓刑考验期满的性质是不执行原判刑罚，而非已经执行完毕。而累犯要求前罪刑罚执行完毕。

3.违反规定。这是指在缓刑考验期内，违反法律法规、监管规定或禁止令，情节严重的，应当撤销缓刑，执行原判刑罚。\footnote{（1）在考验期内违反规定，无论是在考验期内还是在考验期满后被发现，都要撤销缓刑。这一点，与在考验期内犯新罪的处理和理由是相同的。（2）在缓刑考验期内，没有剥夺犯罪分子行使言论、出版、集会、游行等权利。这一点和管制不同。执行管制时，未经执行机关批准，不得行使这些权利。也即，管制期间比缓刑考验期间处罚要严厉一些。}

